% paper_edits_2026-11-12.tex: edits for V_Friction.tex under the awareness route. Pedro applies them. The main thread does not edit the paper.
% Outline behind it: paper_outline_2026-11-12.md.
% BASIS: your save of 11-12 22:11 (Section 7 typed through %D37 without D17 and D18, the seven slips fixed; the new Appendix E typed
%        through %AB2).
% ORDER: PART 1 (apply now, in this order): F4, E2, X1, then C1 and MV1 if you want them.
%        PART 2 (on hold, do not apply yet): P1, H1, G3, H2, S35, T1, F2, M1, M2.
% HOW:   in a SEARCH / REPLACE block, find the SEARCH text (it exists exactly once) and replace it with the REPLACE text.
%        In a RANGE or INSERT block, follow the instruction on its first lines.
% CHANGED 11-12 22:14: rebuilt on your 22:11 save. F3 is gone (you applied all seven). F4 lists two slips in the appendix as
%        typed. E2 is now only what is left of Appendix E, from AB3 on. X1 is unchanged, and you confirmed its centering.
%        T1 is new, on hold, for your comment on the state-vector table.
% RULES in the new text: Oxford comma; \left and \right; tables at [htpb]; short paragraphs; your notes to yourself kept.
%
% YOUR COMMENTS IN THE PAPER, AND WHAT THE FILE DOES WITH EACH
%   D23 "CHECK THIS": untouched. It is your note.
%   D26 (people with one account can open new ones): equation (17) already allows it. Its denominator is the adults with no
%        relationship at j, whether or not they have one elsewhere. Block C1 adds the sentence that says so.
%   D35: you decided to cluster by state. Block S35 (Part 2) rewrites the paragraph and waits for the state-clustered keys.
%   D36 (move to the appendix): block MV1 does it, as a third subsection of Appendix E.
%   D37: you decided to replace the Pix dummy with instrumented use. Block X1 does it.
%   After the state-vector table (Appendix C): "remove the variables that are collected and not in the state vector".
%        The table is generated by make_state_vector_table.py, so IMPR changes it, with your go. My suggestion: drop the block
%        "Collected, not in the state vector" from Panel A and keep Panel B, which is the evidence for leaving three of them out.
%        The other two (Pix use, branch density) are already explained in Section 5. Block T1 (Part 2) is the text that goes
%        with the new table.
%   Demand Residuals (Appendix C): "What is this whole section giving us?" No edit proposed. It is not superseded, for two
%        reasons. Section 6 leans on it for the one threat to the sleepy share that Section 7 does not remove: demand shocks
%        that are serially correlated would be read as carry and push the estimate up, and Panel B is the check. And Panel C
%        (the residual loads on the bank's own branches) is the evidence that motivates branches and advertising in the
%        availability index. Section 7 is not estimated, so it cannot stand in for either. The section could be shorter
%        (AR11 is long), and every number in it changes with the rerun, so I would trim it then.
%   AB2 (the line on cooperatives): left as you have it.

% ######################################################################################################################
% PART 1. APPLY NOW, IN THIS ORDER: F4, E2, X1, and optionally C1 and MV1
% ######################################################################################################################
% ======================================================================================================================
% F4. Two slips in Appendix E as typed. required (the first is in a section title).
% ======================================================================================================================
% SEARCH
at the Bank Boundery}
% REPLACE
at the Bank Boundary}
% SEARCH
more than 500,00 active accounts
% REPLACE
more than 500,000 active accounts
% ======================================================================================================================
% E2. The rest of Appendix E, from AB3 on. required: Section 7 refers to app:shifters:fees, and D37 (in X1) quotes its table.
%     INSERT: after the paragraph that ends with %AB2 and your note, and before \end{document}, add the text between the
%     markers. Skip any paragraph you have typed in the meantime. Counters AB3 to AB12.
% ======================================================================================================================
% ---- BEGIN E2 (insert after the paragraph that ends with %AB2) ----
A comparison of money at institutions on and off Pix within a market is empty. In 2020Q4 the \pn[w]{pix.roster.B.n_mandated} mandated conglomerates held all but R\$\pn[1]{pix.roster.B.sav_nonmandated_brl_mn.2020Q4} million of the R\$\pn[0]{pix.roster.B.sav_total_brl_bn.2020Q4} billion of $\mathrm{B}$ savings, and the $\mathrm{B}$ conglomerates off Pix held at most R\$\pn[1]{pix.roster.B.sav_offpix_brl_mn.max} million of savings in any quarter from 2020Q4 to 2024Q4. %AB3

Participation comes from the BCB's dated lists of Pix participants, which mark each institution as a mandatory or a voluntary participant \parencite{bcb_pix_participantes}, and from the month-end stock of Pix keys by participant \parencite{bcb_pix_chaves}. Participants are matched to conglomerates through the roots of their tax identifiers (CNPJ). %AB4

The second exposure interacts the market's 2019 mobile access with the post-launch indicator. Mobile lines are already a market state in $\boldsymbol{S}_{mt}$, so this exposure identifies a Pix effect that varies with access. It is valid only if 2019 access is unrelated to other post-2020 shocks. Since the emergency aid went mostly to poorer markets, the design controls for Cad\'{U}nico registration interacted with the post-launch indicator. %AB5

\textcite{sarkisyan2024instant} relates the deposits of small relative to large banks in a municipality to local Pix use, and uses high-speed internet access, through 4G/5G mobile networks, as an alternative instrument. The third exposure follows him, with the same outcome and the 2019 4G/5G share, so that our estimates can be compared with his. %AB6

In the first stage of Table \ref{tab:pix_exposure_fs}, usage is the number of distinct individuals who send Pix transfers, per inhabitant, by MCA and quarter from 2020Q4 to 2024Q4. Each exposure is fixed at 2019 levels and scaled by its standard deviation. The regressions include quarter effects, 2019 Cad\'{U}nico registration interacted with quarter effects, and log population, with population weights. %AB7

Despite its small spread, the mandate exposure has a first-stage $F$ statistic of \pn[1]{pix.exposure.mandated.fs.payers.f}. Mobile access gives \pn[1]{pix.exposure.mobile.fs.payers.f} with clusters by MCA and \pn[1]{pix.exposure.mobile.fs.payers.f_state} with clusters by state, and the 4G/5G share \pn[1]{pix.exposure.g45.fs.payers.f}. Together, the three give a joint $F$ of \pn[1]{pix.exposure.fs.joint.payers.f}, and \pn[1]{pix.exposure.fs.joint.payers.f_state} with clusters by state. %AB8

\input{tab_pix_exposure_fs.tex} % Label: tab:pix_exposure_fs

\subsection{Fees and Sweeps}\label{app:shifters:fees}

Account fees are the explicit cost of keeping an account, and they fall mostly on incumbents. In the 2017 to 2018 Household Budget Survey (\textit{Pesquisa de Or\c{c}amentos Familiares}, POF), \pn[p0]{fees.pof.acct.share_paying} percent of families paid an account-keeping fee, R\$\pn[2]{fees.pof.acct.mean_payer} a month on average among those who paid, at January 2018 prices \parencite{ibge_pof2017_2019}. At the five largest banks, the median listed price of the first standardized service package rose from R\$\pn[2]{fees.listed.pkg1.big5.median.2016} in 2016 to R\$\pn[2]{fees.listed.pkg1.big5.median.2024} in 2024, in nominal terms \parencite{bcb_tarifas_mensal}. These are listed ceilings, not the prices clients pay. %AB9

On the public complaints platform \texttt{consumidor.gov.br}, fee disputes make up \pn[p0]{fees.disputes.acct_share.big5} percent of account complaints against the same five banks over 2016 to 2024, and \pn[p0]{fees.disputes.acct_share.digital} percent against Nubank, Inter, and C6 \parencite{senacon_consumidor_dados}. At the five banks the share falls to \pn[p0]{fees.disputes.acct_share.big5.2024} percent in 2024. %AB10

Fees thus make an incumbent account costly to keep and a digital account nearly free. They push balances toward digital banks, so they cannot explain why balances drift toward those banks slowly. A cost that slows the drift would be a hassle or an information cost rather than a fee. %AB11

Moving money within an institution is not frictionless either, since a time deposit redeemed before maturity restarts its tax clock (Appendix \ref{app:clock}) and a new product may require onboarding. Automatic sweeps work the other way, since they move balances within a bank without a decision. Since 2022, Santander's general conditions for new individual accounts make its sweep from checking accounts into certificates of deposit a default feature \parencite{santander_condicoes_pf_2022}, while Bradesco's sweep requires the client to opt in \parencite{bradesco_investfacil_2015,bradesco_investfacil_2020}. %AB12
% ---- END E2 ----
% ======================================================================================================================
% X1. Instrumented Pix use takes the place of the launch indicator. required by your decision of 22:08 (your note on D37).
%     Apply after F3. Six SEARCH / REPLACE pairs: equation (12), two places in D9, D20, the introduction, and the conclusion.
%     Then five paragraphs replaced whole (D7, D27, D28, D29, D37). For each, replace the paragraph that starts with the words
%     shown on the BEGIN line with the text between the markers.
%     Pix use is measured relative to its mean across markets in the quarter (D7, D27). You confirmed this at 22:14.
%     Without it, use is zero before the launch and positive after it, so its coefficient would also carry the national break
%     of 2020Q4. The instruments are the three exposures of 2019, as before.
%     The former D29 (the "dose" as a second comparison) is absorbed: the dose is now the main term.
%     Your note after %D37 is kept.
% ======================================================================================================================
% SEARCH
+\gamma_P\indicator{t \geq 2020\mathrm{Q}4} - \gamma_C
% REPLACE
+ \gamma_P\,\mathrm{Pix}_{mt} - \gamma_C
% SEARCH
\gamma_{P\varnothing}\indicator{t \geq 2020\mathrm{Q}4}
% REPLACE
\gamma_{P\varnothing}\,\mathrm{Pix}_{mt}
% SEARCH
is what the launch did to awareness
% REPLACE
is what Pix did to awareness
% SEARCH
It is read from the part of the response to the launch of Pix that rivals' money shows and new money does not.
% REPLACE
It is read from the part of the response to the use of Pix that rivals' money shows and new money does not.
% SEARCH
The launch of Pix may move both, and the cost is
% REPLACE
The use of Pix may move both, and the cost is
% SEARCH
The launch of Pix may shift both awareness and the cost of moving money
% REPLACE
The use of Pix may shift both awareness and the cost of moving money
% ---- BEGIN X1 D7 (the paragraph that starts: Advertising and branches move awareness. The launch of Pix m) ----
Advertising and branches move awareness. Pix may have moved both frictions. It lowered the cost of moving money across institutions, and it may have raised awareness, since many accounts were opened to receive transfers (Section \ref{sec:inst_setting:taxonomy}). $\mathrm{Pix}_{mt}$ is the market's use of Pix, the number of individuals who send transfers per inhabitant, relative to its mean across markets in the quarter. Its coefficient $\gamma_P$ is therefore the change in availability with the use of Pix, not yet a statement about either friction. The level of the switching cost is not separate from the level of awareness, and both sit in $\gamma_0$. %D7
% ---- END X1 D7 ----
% ---- BEGIN X1 D27 (the paragraph that starts: The launch of Pix in 2020Q4 shifts both indices.) ----
The use of Pix shifts both indices. Since $\mathrm{Pix}_{mt}$ is measured relative to its mean across markets in the quarter, it is zero before 2020Q4 and compares markets after the launch. The national break at the launch is therefore not read as an effect of Pix. It shares its date with the pandemic, the emergency aid, and the start of the tightening cycle, and the demand system of Section \ref{sec:empirical:demand} carries no quarter effects that would absorb them. %D27
% ---- END X1 D27 ----
% ---- BEGIN X1 D28 (the paragraph that starts: The two coefficients separate what the launch did.) ----
The two coefficients separate what Pix did. If it only raised awareness, rivals' money and new money respond alike, $\gamma_{P} = \gamma_{P\varnothing}$. If it only lowered the cost of moving money, $\gamma_{P\varnothing} = 0$. The difference $\gamma_{P} - \gamma_{P\varnothing}$ is the change in the switching cost. Account openings give a second reading of $\gamma_{P\varnothing}$, with two limits. The opening rate \eqref{eq:17} is national, so Pix enters it through the time series alone. And accounts opened to receive transfers reflect the convenience of the account as well as awareness of the institution, so openings may overstate the awareness part. %D28
% ---- END X1 D28 ----
% ---- BEGIN X1 D29 (the paragraph that starts: A second comparison uses the intensity of use,) ----
Usage is an equilibrium object, since attentive markets adopt Pix faster. It is instrumented by three exposures fixed in 2019: the deposit share of the institutions mandated into Pix at launch, mobile access, and the share of fourth- and fifth-generation (4G/5G) mobile coverage (Appendix \ref{app:shifters:pix}). The exposures must be unrelated to the other shocks that followed 2020, market by market. The emergency aid went mostly to poorer markets, so the design controls for Cad\'{U}nico registration in 2019 interacted with quarter effects. %D29
% ---- END X1 D29 ----
% ---- BEGIN X1 D37 (the paragraph that starts: On the launch of Pix, the one estimate in hand is that of Se) ----
On Pix, the one estimate in hand is that of Section \ref{sec:estimation:sleep}, which uses the national date. There the launch did not detectably reduce measured inertia, and the interval is wide. For the comparison across markets, the three exposures predict the use of Pix with a joint first-stage $F$ of \pn[1]{pix.exposure.fs.joint.payers.f_state}, with clusters by state (Appendix \ref{app:shifters:pix}). An equation that reads the cost from net flows alone, source by source, does not converge on the data, so we do not pursue it. %D37 % This is all suggesting to me to abandon the dummy for pix and reaplce it w an instrumented version of the usage.
% ---- END X1 D37 ----
% ======================================================================================================================
% C1. One sentence after equation (17), for your comment on D26. optional.
%     The SEARCH ends at the counter, so your comment after it stays where it is.
% ======================================================================================================================
% SEARCH
the choice probability of $j$ among depositors who consider it. %D26
% REPLACE
the choice probability of $j$ among depositors who consider it. The adults at risk are those with no relationship at $j$, whether or not they hold one elsewhere, so the rate allows for depositors with accounts at several institutions. %D26
% ======================================================================================================================
% MV1. The placebo of the stakes moves to the appendix, for your comment on D36. optional. Apply after E2 (it needs %AB12).
%     Pair 1 is a RANGE of one line: replace the whole paragraph that starts "Part of that strength comes from the lagged deposit
%     weights" (it ends with %D36 and your note) with the text between the markers. Your note is kept after the counter.
%     Pair 2 adds a third subsection at the end of Appendix E. Counter AB13.
% ======================================================================================================================
% ---- BEGIN MV1 pair 1 ----
Part of that strength comes from the lagged deposit weights that any bank attribute would carry. A placebo that shuffles the cost-predicted spreads across institutions shows that the cost shifters carry the rest (Appendix \ref{app:shifters:stakes}). %D36 % eventually this is should be moved to the appendix
% ---- END MV1 pair 1 ----
% SEARCH
while Bradesco's sweep requires the client to opt in \parencite{bradesco_investfacil_2015,bradesco_investfacil_2020}. %AB12
% REPLACE
while Bradesco's sweep requires the client to opt in \parencite{bradesco_investfacil_2015,bradesco_investfacil_2020}. %AB12

\subsection{The Stakes}\label{app:shifters:stakes}

The cost shifters predict the log of the stakes with an $F$ of \pn{pix.stakes.fs.f}. Shuffling the cost-predicted spreads across institutions within the quarter gives a median $F$ of \pn{pix.stakes.fs.f_placebo}, and \pn[p0]{pix.stakes.fs.placebo_share_above} percent of the shuffles reach the real statistic. The lagged deposit weights therefore account for part of the first stage, and the cost shifters carry the rest. %AB13
% ######################################################################################################################
% PART 2. ON HOLD. DO NOT APPLY YET: P1, H1, G3, H2, S35, T1, F2, M1, M2
% ######################################################################################################################
% ======================================================================================================================
% ON HOLD, DO NOT APPLY YET: P1, H1, G3, H2, F2, M1, M2.
%     Why they wait: the numbers in the paper still come from the old market size. The model text should change in the same pass as
%     the estimates, after IMPR's Step 1 and your go on Step 2. The drafts are here so that you can read and correct them now.
%     Still to be written after Step 1: paragraphs E23 to E27 of Section 5 (data-implied shares and the market size) and the
%     market-size appendix. They need the definition of the broad total, the return used for its carry, and the rule that
%     allocates the national pool to markets.
% ======================================================================================================================
% ======================================================================================================================
% P1 (on hold). Section 4, Deposit Dynamics: new money and the awake pool. Replaces paragraphs M18 to M20.
%     Your decision of 11-12: the awake pool is the awake share of the carried wealth plus new money, and the market is liquid
%     financial wealth, not deposits alone. Equations (7-B) and (7-D) keep their tags. What was (1 - phi) M becomes the pool P.
%     Counters: M18 to M23 here, so the paragraph that is M21 today ("Institutions set rates on only two...") becomes M24.
%     The last paragraph makes one claim about the data (deposits grow by more than their interest). Step 1's table must confirm it,
%     and it should then carry a key.
%     RANGE: replace FROM the line that starts   Let $\mathrm{Dep}_{jkmt}$ denote the total deposit stock
%            THROUGH the line that ends          but not directly observed in the data. %M20
% ======================================================================================================================
% ---- BEGIN P1 (on hold) ----
Let $\mathrm{Dep}_{jkmt}$ denote the total deposit stock of type $k$ held at institution $j$ in market $m$ at time $t$. Money that is asleep stays where it was and accrues interest. Money that is awake is placed again. %M18

The awake money of a market has two parts. The first is the awake share of the wealth carried from the previous quarter, $\left(1 - \phi_{mt}\right)C_{mt}$, where $C_{mt}$ is the liquid wealth of the market at the end of the previous quarter, in deposits and in the outside assets, with its accrued return. The second is new money, $N_{mt}$, the wealth that enters the market during the quarter. New money has no account in which to be carried forward, so all of it is awake. The awake pool is $P_{mt} = \left(1 - \phi_{mt}\right)C_{mt} + N_{mt}$. %M19

Deposit balances evolve differently by institution type. For $j \in J_{mt}^{\mathrm{B}}$, the stock of type-$k$ deposits in market $m$ follows %M20
\begin{gather}
	\mathrm{Dep}_{jkmt}^{\mathrm{B}} = \underset{\text{awake: re-choose}}{\underbrace{P_{mt}\,s_{jkmt}^{\mathrm{Act}}(\delta_{mt},\theta_2)}} + \underset{\text{asleep: carry forward}}{\underbrace{\phi_{mt}\left(1 + r_{jkmt-1}^{\mathrm{dep}}\right)\mathrm{Dep}_{jkmt-1}^{\mathrm{B}}}}.
	\label{eq:7-B} \tag{7-B}
\end{gather}
The first term is the awake money that chooses $(j,k)$. The second is the balance carried forward by inactive depositors, with interest at the lagged deposit rate. For $j \in J_t^{\mathrm{D}}$, the institution operates solely at the national level and its deposit stock follows %M21
\begin{gather}
	\mathrm{Dep}_{jkt}^{\mathrm{D}} = P_ts_{jkt}^{\mathrm{Act}}(\delta_t,\theta_2) + \phi_t\left(1 + r_{jkt-1}^{\mathrm{dep}}\right)\mathrm{Dep}_{jkt-1}^{\mathrm{D}},
	\label{eq:7-D} \tag{7-D}
\end{gather}
where $P_t = \sum_{m \in \mathcal{M}}{P_{mt}}$ and $s_{jkt}^{\mathrm{Act}}(\delta_t,\theta_2) = \frac{\sum_{m \in \mathcal{M}}{P_{mt}s_{jkmt}^{\mathrm{Act}}(\delta_{mt},\theta_2)}}{P_t}$ is the pool-weighted average of local choice probabilities, which is well-defined via \eqref{eq:4} but not directly observed in the data. %M22

Without new money the pool is the awake share of a market of given size, $\left(1 - \phi_{mt}\right)M_{mt}$. Deposits grow by more than their interest, and a pool of that form can supply the growth only if the market is many times larger than its deposits. New money supplies it directly. %M23
% ---- END P1 ----
% ======================================================================================================================
% H1 (on hold, rewritten 11-12 20:35). Section 7.1: the law of motion by source, equation (16), D17 and D18.
%     Place: between D16 and D19. It needs P1, since it refers to the pool of (7-B).
%     Change from the earlier draft: the pool with no source is the awake share of the carried wealth outside deposits plus new
%     money. Before, it was the awake share of the market size net of the stocks.
% ======================================================================================================================
% ---- BEGIN H1 (on hold) ----
The balance of $(j,k)$ in a $\mathrm{B}$ market is then %D17
\begin{gather}
	\begin{aligned}
		\mathrm{Dep}_{jkmt}^{\mathrm{B}} = {} & \phi_{mt}\left(1 + r_{jkmt-1}^{\mathrm{dep}}\right)\mathrm{Dep}_{jkmt-1}^{\mathrm{B}} + \left(1 - \phi_{mt}\right)\sum_{j^\prime}{\mathrm{Stock}_{j^\prime mt-1}\,s_{jk\mid j^\prime}} \\
		& + \left[\left(1 - \phi_{mt}\right)C^{\varnothing}_{mt} + N_{mt}\right]s_{jk\mid\varnothing},
	\end{aligned} \label{eq:16}
\end{gather}
where $\mathrm{Stock}_{j^\prime mt-1} = \sum_{k}{\left(1+r_{j^\prime kmt-1}^{\mathrm{dep}}\right)\mathrm{Dep}_{j^\prime kmt-1}}$ is the accrued stock at source $j^\prime$, $C^{\varnothing}_{mt}$ is the carried wealth outside deposits, and $N_{mt}$ is new money. Money that leaves an outside asset has no bank to leave, so it shares the index of new money. This is the installed-base construction of \textcite{nosal2012estimating,shcherbakov2016measuring}, in which last period's stock by source is the population at risk and the conditional share is the transition rate. With $A \equiv 1$ and $\tau_1 = 0$ every source has the share \eqref{eq:4}, the pools add up to $P_{mt}$, and \eqref{eq:16} is \eqref{eq:7-B}. %D18
% ---- END H1 ----
% ======================================================================================================================
% G3 (on hold). The last sentence of D9, which you typed. It describes new money as the market size net of the lagged stocks.
%     The replacement names the allocation of the national pool. Final wording after Step 1 fixes the rule.
% ======================================================================================================================
% SEARCH (on hold)
We do not rely on the level, because new money is measured as the market size net of the lagged stocks and inherits the error of the market size (Appendix \ref{app:estimation:blp:marketsize}). %D9
% REPLACE (on hold)
We do not rely on the level, because the pool with no source is measured for the country and allocated to markets by a rule (Appendix \ref{app:estimation:blp:marketsize}). %D9
% ======================================================================================================================
% H2 (on hold, text unchanged since 19:55). Section 7.3: D38 to D40. Place: after D37.
%     D40 will be rewritten once Step 1 shows the pool weights under the new construction. D38 and D39 stay close to this.
% ======================================================================================================================
% ---- BEGIN H2 (on hold) ----
One discrepancy is now largely accounted for. The estimates of Section \ref{sec:estimation} imply that awake money is \pn[p1]{kappa.k2.world.awake_share} percent of deposits in a quarter, against a median net flow of \pn[p0]{kappa.k0.moving.median.k4} percent of the carried stock for time deposits. The gap sits in two groups of bank, market, and quarter cells. In \pn[p0]{kappa.gap.share.capped} percent of them the estimated sleepy share is above the cap of Appendix \ref{app:estimation:blp:marketsize}. There the model implies a median net flow of \pn[p1]{kappa.gap.median.capped.W0} percent with the estimated share and \pn[p1]{kappa.gap.median.capped.W3} percent with the capped one, against \pn[p1]{kappa.gap.median.capped.data} percent in the data. In another \pn[p0]{kappa.gap.share.truncated} percent the outflow in the data exceeds the awake stock, which the model cannot produce and \eqref{eq:11-B} truncates. %D38

In the remaining \pn[p0]{kappa.gap.share.uncapped} percent of cells the model applies as written, and its median net flow is \pn[p1]{kappa.gap.median.uncapped.W0} percent against \pn[p1]{kappa.gap.median.uncapped.data} percent in the data. Large net flows are not rebooking events either. Changes of branches or of a conglomerate's members are about as frequent among the largest flows as among all cells, with a ratio of \pn[2]{kappa.gap.lift.any}. The discrepancy therefore points to sleepy shares that are too high where they are close to one, and not to flows that are not choices. %D39

The same feature limits what the availability index can deliver. By the current measure of the market size, the outside pool is \pn[p0]{kappa.k0.outside.share_Mhat.median} percent of it at the median market-quarter, so new money is nearly all of the awake pool. Advertising is then read mostly from new money, and rests on its exclusion from utility more than on the comparison of own with rival money. The index of rivals' money carries little weight, so the difference $\gamma_{P} - \gamma_{P\varnothing}$ and the tax-clock term $\gamma_C$ are close to unidentified. The remedy does not lie in the index, and other anchors for the market size do not provide it either. The active balance divided by the awake share is well above the stock on the books when the sleepy share is close to one, whatever the anchor. Both limits therefore trace to the same sleepy shares. %D40
% ---- END H2 ----
% ======================================================================================================================
% S35 (on hold until the keys exist). D35 with clusters by state, your decision of 22:08.
%     The keys of the state-clustered first stage of the stakes are not built yet. The names below are the ones I asked IMPR for.
%     Replace the paragraph that starts "On attention, the lagged Selic coefficient". Your note after %D35 is kept.
%     When these keys exist, the appendix paragraph of MV1 (AB13) should quote the same clustering.
% ======================================================================================================================
% ---- BEGIN S35 (on hold) ----
On attention, the lagged Selic coefficient of Section \ref{sec:estimation:sleep} has the sign that \eqref{eq:15} predicts, since sleepiness falls when the policy rate rises. It is significant under \ref{estimation:single_idx_time} but not under the preferred \ref{estimation:single_idx}, so the current estimates lean toward this channel without establishing it. The stakes of \eqref{eq:14} have a median of R\$\pn[0]{pix.stakes.median} a year. With clusters by state, the cost shifters predict their log with an $F$ of \pn{pix.stakes.fs.f_state}, or \pn{pix.stakes.fs.f_net_state} net of the lagged market structure. The variation across markets within a quarter is weak, so the first stage is revisited before the stakes enter the estimation. %D35 % If this doesn't change after we fix all the bugs, then cluster by market, not states.
% ---- END S35 ----
% ======================================================================================================================
% T1 (on hold until IMPR regenerates the table). Appendix C, paragraph AR5, for your comment after the state-vector table.
%     It goes with a table whose Panel A no longer has the block "Collected, not in the state vector".
% ======================================================================================================================
% SEARCH (on hold)
Panel A gives the role of each state, the share of its variance absorbed by the fixed effects, and its law of motion, and lists the collected candidates left out with the reason. Panel B shows how a second share of a characteristic already in the vector conditions the regressors of \eqref{eq:10}. %AR5
% REPLACE (on hold)
Panel A gives the role of each state, the share of its variance absorbed by the fixed effects, and its law of motion. Panel B shows how a second share of a characteristic already in the vector, collected but left out, conditions the regressors of \eqref{eq:10}. %AR5
% ======================================================================================================================
% NOTE 21:22: you decided that the fee enters the demand model. After the rerun, F2 gives way to text in Section 4 (the fee as a
%     characteristic of checking in equation (1)), Section 5 (how it is built and that it is not instrumented), and a fee row in
%     Tables 5 and 6. I write those blocks when IMPR's plan and numbers are in. Until then F2 is the interim sentence.
% F2 (on hold until the keys exist). One sentence for the footnote of Section 4 on fees (paragraph M1).
%     From IMPR's fees tables of 11-12. You gave the go for the keys at 20:41. Apply once IMPR reports them in paper_numbers_macros.tex.
%     Place: at the end of the footnote that ends "...fees are set in equilibrium to attract clients."
% ======================================================================================================================
% ---- BEGIN F2 (on hold) ----
 Listed prices add little variation across the largest institutions. The standard account package cost between R\$\pn{demand.fee.pkg1.big5.lo.2024} and R\$\pn{demand.fee.pkg1.big5.hi.2024} a month at the five largest in 2024, and a listed price changes in about \pn[p0]{demand.fee.pkg1.chg.rate} percent of quarters.
% ---- END F2 ----
% ======================================================================================================================
% ON HOLD: M1 and M2 (the Findex sentence and the cap sentence of the market-size appendix).
%     You and IMPR concluded today that the market definition itself is wrong (the anchor is an active balance divided by the awake share,
%     not a deposit density) and a comparison of alternatives is running. That subsection will be rewritten once you choose,
%     so these two blocks are withdrawn until then.
% ======================================================================================================================
